\setchapterpreamble[u]{\margintoc}
\chapter{Vectors, Matrices, and Column Space}
\labch{vectors-and-column-space}

\sources{Strang, \emph{Introduction to Linear Algebra}, 6th ed.\ \cite{strang2023ila},
§1.1--1.4 and §2.1; MIT 18.06 \cite{ocw1806}, Lecture~1.}

Linear algebra begins with a bookkeeping question --- how to store a system of
equations --- and almost immediately turns into a geometric one. This chapter
sets up the notation used for the rest of the book and introduces the two
readings of \(A\vect{x}=\vect{b}\) that everything later depends on.

\section{The matrix form of a linear system}
\labsec{matrix-form}

A system of \(m\) linear equations in \(n\) unknowns can be packed into a single
product,
\[
	A\vect{x}=\vect{b},
\]
where \(A\in\R^{m\times n}\) holds the coefficients, \(\vect{x}\in\R^{n}\) holds
the unknowns, and \(\vect{b}\in\R^{m}\) is the target. Throughout this section I
work with
\[
	A=\begin{bmatrix}
		2 & -1 & 0 \\
		-1 & 2 & -1 \\
		0 & -3 & 4
	\end{bmatrix},
	\qquad
	\vect{b}=\begin{bmatrix}0\\-1\\4\end{bmatrix}.
\]

Read one row at a time, the product unpacks into the three equations
\begin{align*}
	2x-y    &= 0,\\
	-x+2y-z &= -1,\\
	-3y+4z  &= 4.
\end{align*}

\marginnote{Nothing has been solved yet. The matrix form is a change of
notation, not a change of content --- but the notation is what makes the two
pictures below visible.}

Substituting \(\vect{x}=\begin{bmatrix}0&0&1\end{bmatrix}\tp\) satisfies all
three equations at once, so this vector solves the system. The interesting
question is not \emph{whether} this particular vector works, but how one would
have found it, and what the two ways of reading the product suggest.

\section{The row picture: equations as constraints}
\labsec{row-picture}

Each row of a three-variable system describes a plane in \(\R^{3}\). A vector
solves the system exactly when it lies on all three planes simultaneously, so
the solution set is the intersection of the three planes. For the example above
that intersection is the single point \((0,0,1)\).

\begin{remark}
The row picture is a picture of \emph{constraints}. Every additional equation
removes a degree of freedom, and a solution is whatever survives all of the
restrictions at once.
\end{remark}

Its usefulness is bounded by the fact that it has to be drawn. With nine
unknowns the same reasoning would require intersecting hyperplanes in
\(\R^{9}\), which no amount of care will make visible. The algebra is
unaffected --- elimination will still run --- but the geometric intuition stops
paying for itself. That is the first reason to look for a second reading of the
same product.

\section{The column picture: targets as combinations}
\labsec{column-picture}

Write \(\vect{a}_1,\vect{a}_2,\vect{a}_3\) for the columns of \(A\). Expanding
the product by columns instead of by rows gives
\[
	A\vect{x}=x\vect{a}_1+y\vect{a}_2+z\vect{a}_3,
\]
so the system asks for coefficients that mix the three columns into the target:
\[
	x\begin{bmatrix}2\\-1\\0\end{bmatrix}
	+y\begin{bmatrix}-1\\2\\-3\end{bmatrix}
	+z\begin{bmatrix}0\\-1\\4\end{bmatrix}
	=\begin{bmatrix}0\\-1\\4\end{bmatrix}.
\]
In this instance the answer can be read off without any calculation: the target
\emph{is} the third column, so \((x,y,z)=(0,0,1)\). The example is convenient,
but the idea it illustrates is not.

\begin{definition}
A vector \(\vect{b}\) is a \emph{linear combination} of
\(\vect{a}_1,\ldots,\vect{a}_n\) when there exist scalars \(x_1,\ldots,x_n\)
with
\[
	\vect{b}=x_1\vect{a}_1+\cdots+x_n\vect{a}_n .
\]
The set of all such combinations is the \emph{column space} of \(A\), written
\(\Col(A)\).
\end{definition}

The definition converts a question about solving into a question about
membership:
\[
	A\vect{x}=\vect{b}\ \text{is solvable}
	\quad\Longleftrightarrow\quad
	\vect{b}\in\Col(A).
\]

\marginnote{The row picture asks whether a candidate satisfies every
restriction. The column picture asks which recipe of the available ingredients
produces the target. Both are correct; they suggest different algorithms.}

Changing the target tests a different recipe built from the same ingredients.
Taking one of the first column and one of the second,
\[
	\begin{bmatrix}2\\-1\\0\end{bmatrix}
	+\begin{bmatrix}-1\\2\\-3\end{bmatrix}
	=\begin{bmatrix}1\\1\\-3\end{bmatrix},
\]
so the system with target \(\begin{bmatrix}1&1&-3\end{bmatrix}\tp\) is solved by
\(\vect{x}=\begin{bmatrix}1&1&0\end{bmatrix}\tp\).

\subsection{When is every target reachable?}
\labsubsec{reachability}

For a square matrix \(A\in\R^{n\times n}\) the natural question is whether
\(A\vect{x}=\vect{b}\) can be solved for \emph{every} \(\vect{b}\in\R^{n}\). In
the language just introduced this asks whether the columns fill the whole space,
\[
	\Col(A)=\R^{n}.
\]

If the columns are linearly independent they span \(\R^{n}\), every target is
reachable, and the coefficients are unique. The matrix is then invertible and
\(\vect{x}=A^{-1}\vect{b}\).

If instead the columns are dependent, at least one of them contributes no new
direction: it already lies in the span of the others. In \(\R^{3}\) the extreme
case is three columns lying in a common plane. Every combination stays in that
plane, so any target off the plane is unreachable, while targets on it are
reached in infinitely many ways.

\begin{remark}
Dependence and invertibility are the same phenomenon seen from two sides.
Independent columns give exactly one solution for each of the \(n\)-dimensional
family of targets; dependent columns give either none or infinitely many,
depending on where the target sits.
\end{remark}

This dichotomy --- invertible against singular --- is the thread that runs
through elimination in \refch{elimination}, the four subspaces in
\refch{four-subspaces}, and the singular value decomposition in \refch{svd}. It
reappears, transformed, whenever a network layer is asked whether it can
represent a given output.

\section{Matrix--vector multiplication, two readings}
\labsec{matvec}

The column viewpoint supplies the definition of the product that is easiest to
reason with. If
\[
	A=\begin{bmatrix}\vert&&\vert\\ \vect{a}_1&\cdots&\vect{a}_n\\ \vert&&\vert\end{bmatrix},
	\qquad
	\vect{x}=\begin{bmatrix}x_1\\\vdots\\x_n\end{bmatrix},
\]
then
\begin{equation}
	A\vect{x}=x_1\vect{a}_1+\cdots+x_n\vect{a}_n .
	\labeq{matvec-columns}
\end{equation}

The rule taught first in school computes one entry at a time: entry \(i\) of
\(A\vect{x}\) is the dot product of row \(i\) of \(A\) with \(\vect{x}\). Both
give the same numbers. Taking
\[
	A=\begin{bmatrix}2&5\\1&3\end{bmatrix},\qquad
	\vect{x}=\begin{bmatrix}1\\2\end{bmatrix},
\]
the column reading gives
\[
	1\begin{bmatrix}2\\1\end{bmatrix}+2\begin{bmatrix}5\\3\end{bmatrix}
	=\begin{bmatrix}12\\7\end{bmatrix},
\]
and the row reading gives the same vector through \(2\cdot1+5\cdot2=12\) and
\(1\cdot1+3\cdot2=7\).

\begin{remark}
The two readings are not rivals; they answer different questions. Row by row is
how the product is \emph{computed} entry by entry. Column by column is what the
product \emph{means}: \(A\vect{x}\) never leaves the column space of \(A\), no
matter which \(\vect{x}\) is supplied.
\end{remark}

That last observation is worth stating on its own, because it is the reason
\refeq{matvec-columns} is the form used throughout this book.

\begin{proposition}
For every \(\vect{x}\in\R^{n}\) the product \(A\vect{x}\) lies in
\(\Col(A)\), and every element of \(\Col(A)\) arises as \(A\vect{x}\) for some
\(\vect{x}\).
\end{proposition}

\begin{proof}
Equation~\refeq{matvec-columns} exhibits \(A\vect{x}\) as a linear combination
of the columns, which is membership in \(\Col(A)\) by definition. Conversely a
combination \(x_1\vect{a}_1+\cdots+x_n\vect{a}_n\) is the product of \(A\) with
the coefficient vector \(\vect{x}=[x_1,\ldots,x_n]\tp\).
\end{proof}

\section{Span, independence, and the shape of a solution set}
\labsec{span-independence}

Two definitions make the earlier discussion precise.

\begin{definition}
The \emph{span} of \(\vect{a}_1,\ldots,\vect{a}_n\) is the set of all their
linear combinations. The vectors are \emph{linearly independent} when
\[
	x_1\vect{a}_1+\cdots+x_n\vect{a}_n=\vect{0}
	\quad\text{forces}\quad
	x_1=\cdots=x_n=0 ,
\]
and \emph{linearly dependent} otherwise.
\end{definition}

\marginnote{Independence is a statement about the zero target only, yet it
controls the solution count for every target at once. That leverage is what
makes it the central definition of the subject.}

Dependence is exactly the statement that some non-zero \(\vect{x}\) satisfies
\(A\vect{x}=\vect{0}\). Suppose such an \(\vect{x}_0\neq\vect{0}\) exists and
\(\vect{x}_p\) solves \(A\vect{x}=\vect{b}\). Then for every scalar \(c\),
\[
	A(\vect{x}_p+c\vect{x}_0)=A\vect{x}_p+cA\vect{x}_0=\vect{b}+\vect{0}=\vect{b},
\]
so one solution has multiplied into infinitely many. This is the whole
explanation for the trichotomy that a linear system obeys.

\begin{theorem}
A system \(A\vect{x}=\vect{b}\) has no solution, exactly one solution, or
infinitely many. It never has, for instance, exactly two.
\end{theorem}

\begin{proof}
If there is no solution the first case holds. If \(\vect{x}_p\) and
\(\vect{x}_q\) are distinct solutions then \(\vect{x}_0=\vect{x}_p-\vect{x}_q\)
is non-zero and satisfies \(A\vect{x}_0=\vect{0}\); the displayed calculation
above then produces a distinct solution for every scalar \(c\), so there are
infinitely many. Hence two or more solutions already forces infinitely many.
\end{proof}

The set of vectors killed by \(A\) therefore deserves a name of its own; it is
the nullspace, and \refch{four-subspaces} is largely about it.

\section{Notation used in this book}
\labsec{notation}

Lower-case bold letters denote column vectors, capitals denote matrices, and
\(A\tp\) is the transpose. The columns of \(A\) are \(\vect{a}_1,\ldots\); the
entry in row \(i\) and column \(j\) is \(a_{ij}\). The symbol \(\R^{n}\) is the
set of real column vectors of length \(n\), and \(I\) is the identity matrix
whose size is always clear from context.

Vectors are columns unless stated otherwise, so a row vector is written
\(\vect{v}\tp\). Data matrices are the one place where this convention grates:
a batch of \(n\) examples with \(d\) features is stored as \(X\in\R^{n\times d}\)
with one example per \emph{row}, which is the convention every numerical library
uses. Whenever both appear together the shapes are stated explicitly rather than
left to be inferred.
